\documentclass[11pt]{article}
\usepackage[a4paper,margin=1in]{geometry}
\usepackage{fontspec}
\usepackage[boldfont=false]{xeCJK}
\setCJKmainfont{FandolSong-Regular.otf}
\usepackage{amsmath}
\usepackage{amssymb}
\usepackage{array}
\usepackage{booktabs}
\usepackage{graphicx}
\usepackage{adjustbox}
\usepackage{enumitem}
\setlist[itemize]{topsep=2pt,partopsep=0pt,parsep=0pt,itemsep=2pt,leftmargin=1.4em}
\setlist[enumerate]{topsep=2pt,partopsep=0pt,parsep=0pt,itemsep=2pt,leftmargin=1.6em}
\usepackage{parskip}
\usepackage{fvextra}
\fvset{breaklines=true,breakanywhere=true,fontsize=\small}
\setlength{\parindent}{0pt}

\begin{document}

\begin{center}
  {\LARGE\bfseries Hydroxy compounds}\\[0.35em]
  {\large A Level Chemistry}
\end{center}
\vspace{0.6em}

\section*{Alcohols with acyl chlorides}

At A Level you meet one more reaction of an \textbf{alcohol} 醇: it reacts with an \textbf{acyl chloride} 酰氯 to make an \textbf{ester} 酯. This works faster and more completely than using a carboxylic acid. For example, ethanol and ethanoyl chloride give ethyl ethanoate, plus fumes of $\text{HCl}$.

\section*{Phenol}

\textbf{Phenol} 苯酚 has an $\text{–OH}$ group joined directly to a \textbf{benzene} 苯 ring. This changes the chemistry of both the $\text{–OH}$ group and the ring.

\subsection*{Making phenol}

Cool \textbf{phenylamine} 苯胺 with $\text{NaNO}_2$ and dilute acid below $10\,°\text{C}$ to make a \textbf{diazonium salt} 重氮盐. Warming this salt with water then gives phenol (and nitrogen gas).

\subsection*{Reactions of phenol}

\begin{itemize}
\item with a base such as $\text{NaOH(aq)}$: phenol reacts to give sodium phenoxide and water. (Ordinary alcohols do \textbf{not} react with $\text{NaOH}$ — this shows phenol is more acidic.)

\item with sodium metal: gives sodium phenoxide and hydrogen.

\item with a diazonium salt in $\text{NaOH(aq)}$: forms a coloured \textbf{azo compound} 偶氮化合物 (used in dyes).

\item \textbf{nitration} 硝化 with \textbf{dilute} $\text{HNO}_3$ at room temperature: gives a mixture of 2-nitrophenol and 4-nitrophenol.

\item \textbf{bromination} 溴化 with bromine water at room temperature (no catalyst): gives a white precipitate of 2,4,6-tribromophenol.

\end{itemize}

\subsection*{Acidity of phenol}

Phenol loses its $\text{H}^+$ to form a phenoxide ion. This ion is stabilised because its negative charge is spread (delocalised) into the ring. So phenol's \textbf{acidity} 酸性 is higher than that of water or ethanol:

\[
\text{ethanol} < \text{water} < \text{phenol}
\]

(Phenol is still a weak acid — weaker than a carboxylic acid.)

\subsection*{Why the conditions are milder than for benzene}

A lone pair on the phenol oxygen is partly \textbf{delocalised} 离域 into the ring. This makes the ring more electron-rich, so it attracts electrophiles more strongly. Phenol therefore reacts faster and under much milder conditions than benzene — no catalyst is needed, and dilute reagents work at room temperature.

The $\text{–OH}$ group directs new substituents to the 2-, 4- and 6-positions. The same ideas apply to other phenolic compounds, such as naphthol.


\end{document}
